% ============================================================
% Capítulo: Amenazas a la Validez (Bloque B.12, D5)
% ============================================================
\chapter{Amenazas a la validez}
\label{cap:amenazas-validez}

Este capítulo analiza, en las 4 categorías estándar de amenazas a la
validez de un estudio empírico de ingeniería de software, las
limitaciones \textbf{reales} ya documentadas en el repositorio -- no una
lista genérica de amenazas de relleno que no correspondan a algo
efectivamente observado o declarado durante este proyecto.

\section{Validez de constructo}

La amenaza de constructo más concreta de este proyecto está en el propio
diseño del experimento de rendimiento (\autoref{sec:cache-redis},
\texttt{docs/mediciones/perf/REPORT.md}): el escenario
\texttt{cache\_frio} se diseñó para medir el costo de una consulta al
catálogo \emph{sin} caché (el constructo de interés, para contrastarlo
contra \texttt{cache\_caliente}), pero su implementación real -- un PRNG
que elige una página aleatoria entre 0 y 999 en cada petición -- hace que
la gran mayoría de esas páginas queden fuera de rango de la tabla
\texttt{libro} (poblada con datos de ejemplo de desarrollo, no con miles
de registros), de modo que Spring Data devuelve una \texttt{Page} vacía
sin necesidad de traer filas reales. En la práctica,
\texttt{cache\_frio} mide sobre todo el costo de una consulta
\texttt{COUNT}~+~\texttt{SELECT} con resultado vacío, \textbf{no} el
costo real de traer y serializar una página completa de libros sin
caché -- una construcción operacional distinta de la que el nombre del
escenario sugiere. La comparación ``caché caliente vs.\ caché frío'' de
este experimento, por tanto, \textbf{no aísla limpiamente el efecto del
caché}: parte de la diferencia de latencia observada entre ambos
escenarios refleja también el costo distinto de las consultas que cada
uno ejecuta, no únicamente la presencia o ausencia de caché. Esta
limitación está documentada explícitamente en el propio
\texttt{REPORT.md} (punto~3 de su ``Análisis breve''), no descubierta
recién en este capítulo -- se reproduce aquí porque es exactamente el
tipo de amenaza de constructo que este capítulo debe declarar, en vez de
reemplazarla por una amenaza genérica que no corresponda al proyecto
real.

\section{Validez interna}

La corrida~1 de las 5 corridas de k6 (\autoref{cap:diseno-arquitectura},
\texttt{docs/mediciones/perf/REPORT.md}) presenta una cola de latencia
anómala en \texttt{cache\_caliente} (p95$=$107{,}31\,ms, frente a un
rango de 6{,}25--12{,}36\,ms en las 4 corridas siguientes) que no se
repite en ninguna corrida posterior. La explicación documentada es
\emph{warm-up} de JVM/JIT y del \emph{pool} de conexiones de HikariCP:
la corrida~1 fue la primera ejecución de carga real contra un contenedor
\texttt{sgb\_backend} que llevaba varias horas sin tráfico significativo.
Esto constituye una variable extraña (\emph{confound}) no controlada
directamente por el diseño experimental -- no se descartó la corrida ni
se repitió hasta obtener un resultado más limpio (lo que habría sido, en
sí mismo, una amenaza distinta: sesgo de selección de datos), sino que
se documentó y se incluyó en el agregado de las 5 corridas, confirmando
además mediante el test estadístico pareado que el patrón
\texttt{caliente}~$>$~\texttt{frio} se sostiene igual si se considera
solo el subconjunto de las corridas 2--5. Una amenaza de validez interna
distinta, propia de la auditoría de seguridad
(\autoref{cap:diseno-arquitectura}), es que el mismo equipo que
implementó el sistema es quien ejecutó y documentó la auditoría OWASP: no
existe una segunda parte independiente que haya intentado explotar el
sistema sin conocimiento previo de su código fuente, lo que puede
subestimar la superficie real de vulnerabilidades frente a una prueba de
penetración por un equipo externo.

\section{Validez externa}
\label{sec:validez-externa}

La amenaza de validez externa más relevante y todavía sin resolver es el
tamaño de muestra de la evidencia de usabilidad (System Usability Scale,
SUS): la guía de esta entrega exige un mínimo de $N\geq15$ participantes,
y esa evidencia sigue \textbf{pendiente} al momento de escribir este
capítulo (\texttt{OBS-08} en
\texttt{docs/observaciones/OBSERVACIONES.md}, ``en progreso, asignada a
Panama'', sin commit real todavía). El protocolo completo -- plantilla de
consentimiento informado, tarea de \emph{onboarding} de dos pasos,
código de participante anónimo -- ya está definido y validado, pero no
ejecutado; este capítulo no puede, por tanto, evaluar la validez externa
del resultado SUS en sí (todavía no existe), solo señalar el riesgo de
tamaño de muestra que enfrentará cuando se ejecute. Ese riesgo no es
menor: \citet{baltes2022sampling} documentan, en su revisión crítica de
prácticas de muestreo en investigación de ingeniería de software, que
muestras pequeñas y no aleatorizadas -- exactamente el perfil esperable
de un estudio SUS de PFC con participantes reclutados por conveniencia
dentro de la comunidad universitaria -- limitan severamente la
generalización de cualquier conclusión más allá de la muestra concreta
estudiada, y recomiendan reportar explícitamente el método de
reclutamiento y sus límites en vez de presentar el resultado como
representativo de una población más amplia sin justificarlo. Cuando la
muestra SUS se ejecute, este capítulo deberá actualizarse citando el
tamaño real obtenido, el método de reclutamiento, y aplicando ese mismo
criterio de \citet{baltes2022sampling} a la interpretación del puntaje.
Una segunda amenaza de validez externa, más estructural: todos los
hallazgos de este proyecto provienen de un único contexto (una biblioteca
universitaria ecuatoriana de bajo volumen, con un equipo de 3 personas
sin dedicación exclusiva) -- ninguna conclusión aquí debería
generalizarse sin más a bibliotecas de otro tipo (públicas de gran
escala, especializadas, con perfiles de carga o de usuario distintos).

Una tercera amenaza de validez externa, distinta de las dos anteriores,
es el alcance del propio acervo bibliográfico reunido para la revisión de
trabajos relacionados (\autoref{cap:trabajos-relacionados}): de las 33
referencias que lo componen -- cada una verificada individualmente contra
el registro oficial de Crossref (autor, DOI y venue reales; ese mismo
proceso de verificación detectó y corrigió 5 errores de atribución de
fuente durante su construcción, evidencia del rigor con que se aplicó) --
solo \textbf{12 califican como de ``alto impacto''} en el sentido estricto
que exige el criterio~D6 de la guía (revista con cuartil JCR Q1/Q2, o
conferencia CORE~A del tipo ICSE/FSE/ASE/MSR/EASE/ESEM), por debajo del
umbral de 20 que exige el nivel ``Excelente'' de la rúbrica -- aunque sí
se cumple el umbral de 30 referencias totales. Esto no refleja un déficit
de esfuerzo de búsqueda: el dominio concreto de este PFC (gestión
bibliotecaria universitaria con IA aplicada) es un nicho de aplicación con
un volumen de publicación considerablemente menor en los venues de primer
nivel de Ingeniería de Software que otras áreas más centrales de la
disciplina (pruebas de software, arquitectura, DevOps); la literatura
realmente disponible sobre este dominio se concentra, en cambio, en
revistas especializadas de ciencias de la información (típicamente Q2) y
en conferencias regionales o temáticas -- reales e indexadas, pero fuera
del conjunto CORE~A que domina la producción de Ingeniería de Software
general. Frente a esta limitación real y ya verificada de forma
exhaustiva, la decisión tomada fue reportar el número exacto obtenido (12
de 33) en vez de completar el conteo con citas de venues de calidad no
verificable -- el mismo criterio de verificación cruzada contra Crossref
que corrigió los 5 errores de atribución mencionados arriba es, en sí
mismo, la evidencia de que ese criterio se aplicó de forma consistente en
todo el acervo, no solo donde convenía a un conteo más favorable.

\section{Validez de conclusión}

La comparación estadística pareada entre \texttt{cache\_caliente} y
\texttt{cache\_frio} (Wilcoxon de rangos con signo sobre el p95 de cada
una de las 5 corridas, \texttt{scripts/perf-analysis.py}) es el ejemplo
más concreto de amenaza a la validez de conclusión de este proyecto: el
estadístico exacto de dos colas con $n=5$ pares tiene un valor mínimo
alcanzable de exactamente $p=0{,}0625$ cuando las 5 diferencias apuntan
en la misma dirección sin excepción -- que es precisamente lo que
ocurre aquí (\texttt{cache\_caliente} $>$ \texttt{cache\_frio} en las 5
corridas) --, por lo que el resultado \textbf{no alcanza el umbral
convencional de significancia} ($p<0{,}05$) pese a un tamaño de efecto
grande y consistente (Cliff's delta $=-0{,}68$). Tratar este resultado
como ``no hay diferencia'' sería una conclusión estadísticamente
incorrecta: el límite real es la potencia estadística del tamaño de
muestra mínimo exigido por la guía ($n=5$ corridas), no evidencia
sustantiva en contra del efecto observado -- más corridas subirían la
potencia sin que se espere que cambie la dirección de la conclusión, ya
consistente en las 5 disponibles. Esta interpretación se documenta de
forma explícita en \texttt{docs/mediciones/perf/REPORT.md} en vez de
reportar solo el $p$-valor sin su contexto de potencia, que por sí solo
induciría a la conclusión errónea de ausencia de efecto. Una segunda
amenaza de conclusión, ya señalada en la sección de validez de
constructo de este mismo capítulo, es que ese mismo resultado
($p=0{,}0625$, efecto grande) se interpreta sobre una comparación que no
aísla limpiamente el constructo de interés (el efecto del caché en sí) --
ambas amenazas se refuerzan mutuamente y deben leerse juntas, no de
forma aislada, para no sobre-interpretar la magnitud exacta del efecto
del caché reportada aquí.
